\scanpage{6}
\uline{Conséquence.} Soit $k_0=\mathcal A_0$, c'est un sous-corps de $k$. Soit $L(k_0)$ la catégorie des alg. de Lie de dim finie sur $k_0$. Par Hausdorff on a un foncteur
\[
L(k_0)\xrightarrow[\text{Hausdorff}]{}\operatorname{Grf}(k_0)\xrightarrow[\text{ext. de la base}]{}\operatorname{Grf}(k)
\]
où ``Grf'' désigne la \uline{catégorie des groupes formels}.
Ceci dit, ce foncteur se factorise canoniquement par les $\mathcal A$-groupes (formels) sur $k$. De façon plus précise, soit $\overline W(\unsure{V_0|\mathfrak g_0})$ le complété à l'origine \deleted{d'un}\inserted{de $V_0\otimes_{k_0}k$}\unsure{d'un} $k_0$-espace vectoriel\inserted{de dim finie} \deleted{\unsure{$V_0$}}. C'est un $\mathcal A$-schéma formel.
Les formules de Hausdorff, si $V_0$ est une Alg de Lie, définissent une loi de groupe
\[
\overline W(V_0)\times\overline W(V_0)\unsure{-|\to}\overline W(V_0)
\]
\uline{Cette dernière est une $\mathcal A$-loi}; \deleted{Pour ceci, il suffit de voir que le système de Pfaff associé est défini par une $\mathcal A$-loi\deleted{\inserted{\unsure{$\mathcal A$-forme}}}} et si
$V_0\unsure{-|\to}\unsure{\mathfrak h_0|h_0}$ est un hom d'algèbres de Lie sur $k_0$, l'hom de groupes formels $\overline W(V_0)\unsure{-|\to}\overline W(\mathfrak h_0)$ est un $\mathcal A$-hom.
Utilisant le th. d'\unsure{Ado|Abel}, on est ramené à prouver que si $V_0\subset\unsure{\mathfrak{gl}(n,k_0)|\mathfrak{gl}(n,k_0)_a}$, alors le sous-groupe formel de $\operatorname{Gl}(n,k_0)$ qu'il définit est un sous-$\mathcal A$-schéma formel $\ldots$
\margin{$-$}
\margin{5}
